\documentclass[11pt,a4paper]{article}
\usepackage[utf8]{inputenc}
\usepackage[T1]{fontenc}
\usepackage{lmodern}
\usepackage[french]{babel}
\usepackage{amsmath,amssymb,mathtools}
\usepackage{icomma}
\usepackage{xcolor}
\usepackage[most]{tcolorbox}
\usepackage{enumitem}
\usepackage[a4paper,margin=2.2cm,headheight=26pt,headsep=10pt]{geometry}
\usepackage{fancyhdr}
\usepackage[hidelinks]{hyperref}
\definecolor{primary}{RGB}{222,49,99}
\definecolor{primarydark}{RGB}{150,22,55}
\tcbset{boxbase/.style={enhanced,breakable,boxrule=0.9pt,arc=2.5pt,
  left=3mm,right=3mm,top=2.3mm,bottom=2.3mm,fonttitle=\bfseries,coltitle=white}}
\newtcolorbox[auto counter]{exo}[1][]{boxbase,colback=primary!4,colframe=primary,
  title={\textsc{Exercice}~\thetcbcounter\ifx\relax#1\relax\relax\else\textmd{ --- \itshape #1}\fi}}
\newcommand{\R}{\mathbb{R}}
\newcommand{\ds}{\displaystyle}
\pagestyle{fancy}\fancyhf{}
\fancyhead[L]{\small\textcolor{primarydark}{\textbf{Thème 2} — Logarithmes}}
\fancyhead[R]{\small\textcolor{primarydark}{\textsc{Moyen}}}
\fancyfoot[C]{\small\thepage}
\renewcommand{\headrulewidth}{0.8pt}
\begin{document}

\begin{center}
{\LARGE\bfseries\color{primarydark} Niveau Moyen — Exercices}\\[2pt]
{\color{primary}\rule{0.45\textwidth}{1.2pt}}\\[3pt]
{\small\itshape On enchaîne les propriétés. Aucune calculatrice.}
\end{center}
\vspace{1mm}

\begin{exo}[condenser en un seul logarithme]
Écris chaque expression comme un \emph{seul} logarithme (avec $x,y>0$) :
\begin{enumerate}[label=\alph*),leftmargin=2em,topsep=1pt,itemsep=1pt]
  \item $\ln x+\ln y$ \qquad
  \item $2\log_a x-\log_a y$ \qquad
  \item $\ds\frac12\ln x$ \qquad
  \item $\log 5+\log 2$ \qquad
  \item $3\ln x-\ln x$
\end{enumerate}
\end{exo}

\begin{exo}[simplifier des expressions logarithmiques]
Simplifie au maximum :
\begin{enumerate}[label=\alph*),leftmargin=2em,topsep=1pt,itemsep=1pt]
  \item $\ln e^{2}+\ln 1$ \qquad
  \item $\ds\ln\sqrt{e}$ \qquad
  \item $\log_{10}10^{2a}$ \qquad
  \item $\ds\ln e^{3}-\ln\frac1e$ \qquad
  \item $\log_a a^{5}$
\end{enumerate}
\end{exo}

\begin{exo}[repérer un carré parfait dans un logarithme]
Simplifie en factorisant d'abord l'argument (pense aux produits remarquables), $x\in\R$ :
\begin{enumerate}[label=\alph*),leftmargin=2em,topsep=1pt,itemsep=1pt]
  \item $\ln\bigl(e^{2x}+2e^{x}+1\bigr)$ \qquad
  \item $\ln\bigl(x^{2}+6x+9\bigr)$ (avec $x>-3$)
  \item $\ds\log_{10}\bigl(100\,x^{2}\bigr)$ (avec $x>0$)
\end{enumerate}
\end{exo}

\begin{exo}[équations où l'inconnue est en exposant]
Résous dans $\R$ (exprime la solution avec $\ln$ ou $\log$) :
\begin{enumerate}[label=\alph*),leftmargin=2em,topsep=1pt,itemsep=1pt]
  \item $e^{2x}=10$ \qquad
  \item $2\times 10^{5x}=5$ \qquad
  \item $3^{x}=7$ \qquad
  \item $e^{x+1}=e^{2}\cdot e^{x}$
\end{enumerate}
\end{exo}

\begin{exo}[inéquations logarithmiques : le domaine décide]
Pour chacune, écris d'abord la condition d'existence (C.E.), puis résous :
\begin{enumerate}[label=\alph*),leftmargin=2em,topsep=1pt,itemsep=1pt]
  \item $\ln(x-1)\geq\ln 2$
  \item $\ln(2x)\leq\ln 6$
  \item $\log_{10}(x+3)>1$
\end{enumerate}
\end{exo}

\end{document}
